\section{Approach to the Three Threats}
\label{sec:threats}

Three threats were chosen from the syllabus: \emph{massive denial-of-service disruptions},
\emph{junk email} and \emph{route and signal hijacking} (as DNS hijacking).
Table~\ref{tab:threats} summarises detection, response and the measured result on the main (k3s)
platform. Each proof of concept is an automated script that prints a threat card, runs the attack
from separate containers, checks detection, response, false positives and recovery, and stores the
raw evidence (alert JSON, firewall rules, logs, Prometheus queries). Every test passed on both
platforms: DoS 17/17 checks, junk email 18/18, DNS 10/10.

\begin{table}[ht]
\centering\small
\caption{Threats, detection, action-response and result on k3s.}
\label{tab:threats}
\begin{tabularx}{\textwidth}{@{}l Y Y Y@{}}
\toprule
Threat & Detection & Action-response & Result (k3s) \\
\midrule
HTTP flood (DoS) & Wazuh \rid{100210}: 15 requests in 10~s from one source; Prometheus rate, p95, in-flight alerts & \code{k3s-block}: drop the source in \code{raw/PREROUTING} on the node, 60~s & alert 2.2~s, blocked 2.8/2.9~s after start; legit client 91\%$\to$100\% \\
Junk email & \rid{100220} (spam), \rid{100223} (quarantined), \rid{100222}: 10 mails in 30~s from one source & spam held by Postfix; flood source \code{firewall-drop} 60~s; relay refused & spam not delivered; flood blocked 2.8~s; no false positive \\
DNS hijacking & FIM real time: \rid{550}$\to$\rid{100230} on the zone file & \code{restore-dns}: restore golden zone, restart dnsmasq & alert 0.1~s, restored 0.2~s; 0~s exposure seen \\
\bottomrule
\end{tabularx}
\end{table}

% ------------------------------------------------------------------------------------------
\subsection{Massive denial of service (HTTP flood)}

\paragraph{Description.}
An attacker sends far more requests than a service can handle, so legitimate requests are delayed
or dropped. At large scale the sources are many (botnets such as Mirai~\cite{antonakakis2017mirai});
the PoC uses a few sources, which is the case the defence targets.

\paragraph{Why it is relevant here.}
Analytics is the most exposed and most expensive endpoint: one \code{/stats?limit=1000} call makes
analytics query the registry and then read 1000 rows from storage. A flood on analytics is
therefore amplified onto storage, the single writer of the whole platform, and degrades ingestion
in every location.

\paragraph{How it can be exploited (PoC).}
Two attacker containers each run 20 concurrent workers against the dc1 analytics NodePort for
25~s, while a legitimate client sends about one request per second, below the detection threshold.
The test has four phases: (1) baseline, legitimate client alone; (2) unmitigated attack, with the
node's Wazuh agent stopped; (3) the same attack with detection and response; (4) recovery after the
block expires.

\paragraph{Monitoring and mitigation metrics.}
Every request produces a JSON access-log line that the node agent forwards (rule \rid{100200}).
Rule \rid{100210} fires when 15 such events from the same \code{srcip} arrive within 10~s.
Sources that are the platform's own services are excluded twice: Incus service addresses through
a CDB list, and the Kubernetes pod network through rule \rid{100209}, a level-0 child of
\rid{100200} that ends pod-to-pod requests before \rid{100210} counts them. On the Prometheus side,
\code{SSLEHighRequestRate}, \code{SSLEHighP95Latency} and \code{SSLEHighConcurrentRequests} watch
rate, latency and backlog per job, and \code{SSLERootCauseRequestFlood} combines them.

\paragraph{Action-response.}
Rule \rid{100210} triggers an active response on the agent that saw the flood, for 60~s. The stock
\code{firewall-drop} adds a DROP rule to the \code{INPUT} chain, which works for the Incus
containers but not for Kubernetes: NodePort traffic is translated by \code{kube-proxy} and then
\emph{forwarded} to the pod, so it never crosses \code{INPUT}. A custom script, \code{k3s-block},
implements the Wazuh active-response protocol and inserts the DROP into the \code{raw} table's
\code{PREROUTING} chain, before \code{kube-proxy}. Because of \code{externalTrafficPolicy: Local},
the node that sees the flood is the node the attacker uses, so blocking locally is enough.

\paragraph{Post-mortem.}
Table~\ref{tab:ddos} shows the legitimate client in each phase. On k3s the alert came 2.2~s after
the attack started and both attackers were blocked within 2.9~s; the legitimate client was never
flagged or blocked, both attackers were unblocked after 60~s, and Prometheus raised
\code{SSLEHighRequestRate} during the unmitigated phase.

\begin{table}[ht]
\centering\small
\caption{DoS test: legitimate client success rate and latency (p50 / p95, ms).}
\label{tab:ddos}
\begin{tabular}{@{}l cc cc@{}}
\toprule
 & \multicolumn{2}{c}{k3s (main)} & \multicolumn{2}{c}{Incus (baseline)} \\
\cmidrule(lr){2-3}\cmidrule(lr){4-5}
Phase & success & p50 / p95 & success & p50 / p95 \\
\midrule
1 baseline    & 100\% & 13 / 32     & 100\% & 11 / 14 \\
2 unmitigated & 91\%  & 900 / 1890  & 0\%   & 2003 / 2003 \\
3 mitigated   & 100\% & 12 / 22     & 77\%  & 10 / 2002 \\
\midrule
first \rid{100210} alert & \multicolumn{2}{c}{2.2~s} & \multicolumn{2}{c}{12.0~s} \\
attackers blocked after  & \multicolumn{2}{c}{2.8 / 2.9~s} & \multicolumn{2}{c}{13.8 / 12.4~s} \\
\bottomrule
\end{tabular}
\end{table}

What failed along the way is as instructive as the final result:
\begin{itemize}
  \item \textbf{Self-inflicted outage (Incus).} In the first automated run, storage and the registry
        saw analytics' fan-out calls as a flood and blocked analytics itself, so the defence took
        the service down. Internal sources are now excluded. The trade-off: a compromised internal
        service flooding another one is not blocked.
  \item \textbf{The same outage again on k3s.} Pod addresses change on every restart, so they
        cannot be listed. A \code{10.42.} prefix entry in the CDB list never matched, and a negated
        CIDR condition in the rule behaved as a positive match, so the analytics pod was blocked
        again. Rule \rid{100209} fixed it and keeps internal HTTP~5xx errors visible for root-cause
        correlation.
  \item \textbf{New pods were invisible to the SIEM.} Every new pod writes a new log file; the agent
        looked for new files only every 64~s, so the first minute of a new pod's traffic was never
        analysed. It now scans every 5~s.
  \item \textbf{Probes amplified overload.} With a 1~s liveness timeout, the kubelet killed pods
        that were slow but healthy (exit 137), which during a flood would restart the attacked pod.
        Liveness now tolerates about 60~s of silence; readiness still removes a struggling pod from
        the Service.
  \item \textbf{Symptoms differ by platform.} On Incus the flood made requests \emph{fail}; on k3s
        the pod limits kept the service answering, slowly (p50 900~ms). The check of the test was
        widened to accept a lower latency as evidence of mitigation.
  \item \textbf{An overloaded service delays its own detection.} Detection took 12.0~s on Incus and
        2.2~s on k3s, but the SIEM pipeline was fast on both: the first \rid{100210} alert came
        0.3~s (Incus) and 1.1~s (k3s) after the fifteenth logged request. The difference is that a
        request is logged when its response is sent. The saturated Incus analytics answered in
        about 2.3~s, so the 15 log lines the rule needs only existed about 9~s into the attack; the
        k3s pod answered in about 0.6~s.
\end{itemize}

Table~\ref{tab:ddos} shows single runs. Repeating the k3s test five times
(Section~\ref{sec:repeat}) gave the same picture in three runs, but in two the analytics pod was so
overloaded that detection took 17.4~s or did not happen.

\paragraph{Distributed flood (false negative).}
A per-source threshold cannot see an attack spread over many sources. This was tested with 40
source addresses (secondary addresses on the two attacker containers), each sending one request
per second, below the 1.5 per second that \rid{100210} needs, for 40~s against the same
endpoint; the generator skips missed slots instead of sending them in bursts. The service received
about 40~req/s in total (as measured by Prometheus) and Wazuh raised \emph{no} \rid{100210} alert
and added no block, while Prometheus fired \code{SSLEHighRequestRate} and
\code{SSLERootCauseRequestFlood} on analytics and storage. At this rate the legitimate client was
not seriously affected (p50 15~ms before, 13~ms during; p95 46 to 226~ms), but the attack was only
visible to Prometheus, which has no automatic response. A botnet would scale the same pattern.

% ------------------------------------------------------------------------------------------
\subsection{Junk email}

\paragraph{Description.}
Unsolicited bulk mail: unwanted content (advertising, phishing) and high-volume sending that
consumes the mail server's resources, plus abuse of the server as an open relay to send spam to
third parties.

\paragraph{Why it is relevant here.}
The mail server is the operators' channel (alerts, notifications) and is reachable from the lab
network. A flood fills its queues; an open relay would make the platform a spam source and get it
blacklisted; phishing reaching operators is the usual first step of an intrusion.

\paragraph{How it can be exploited (PoC).}
From an attacker container the test (1) sends a legitimate message from the legitimate client,
(2) sends a message containing the GTUBE string~\cite{spamassassin_gtube}, (3) tries to relay a
message to an outside domain, (4) opens 12 SMTP sessions in parallel, and (5) checks recovery after
the block.

\paragraph{Monitoring and mitigation metrics.}
Postfix passes every message through SpamAssassin (threshold 5.0; Bayes is enabled in the
configuration but has never been trained, so scores come from the rules alone). Rule
\rid{100220} matches \emph{spamd: identified spam}; \code{header\_checks} moves any message with
\code{X-Spam-Flag: YES} to the hold queue, and Postfix's log line becomes rule \rid{100223}. Rule
\rid{100221} marks each accepted message with its client address, and \rid{100222} fires on 10 of
them from the same source within 30~s. Relaying is refused for any client outside
\code{mynetworks} (localhost only).

\paragraph{Action-response.}
Two layers: the \emph{mail server} itself quarantines spam (hold queue, released by an operator)
and refuses relaying; \emph{Wazuh} responds to a flood (\rid{100222}) with \code{firewall-drop} of the
source on the mail container for 60~s. The spam alerts \rid{100220}/\rid{100223} are alert-only.

\paragraph{Post-mortem.}
On both platforms the legitimate message was delivered and scored clean; the GTUBE message was
flagged, quarantined and not delivered; the relay attempt was refused (\emph{Relay access
denied}). For the flood, the first \rid{100222} alert came 2.2~s after it started and the source was
blocked 2.8~s after it started on the k3s run (1.7~s and 2.0~s on the Incus run). The legitimate
sender was never blocked and could send during the block, and the attacker could connect again after
60~s. How many of the 12 flood messages got through is a race between the parallel sessions and the
block: 2 completed in the Incus run and all 12 in the k3s run. With a threshold of 10 messages, up
to about 10 messages from a new source are accepted before the block.

GTUBE only proves that the pipeline works. To measure detection, 500 spam and 500
legitimate (\emph{easy ham}) messages from the SpamAssassin public corpus (2003) were passed through the
production \code{spamd} (Table~\ref{tab:corpus}); Bayes auto-learning left no database behind. At the
configured threshold 70\% of the spam was caught and 10\% of the legitimate mail would have been
quarantined; lowering the threshold to 3.0 catches more spam but flags almost half of the
legitimate mail. The corpus is old, so its spam may miss current rules and its legitimate mail may
trip date-based ones; which rules caused the false positives was not analysed. Without Bayes
training and network tests, content filtering alone is weak, which is why the quarantine is
reviewed by an operator rather than deleting mail. Other limitations: a sender that rotates
addresses or stays under 10 messages per 30~s is not blocked.

\begin{table}[ht]
\centering\small
\caption{SpamAssassin on 500 spam + 500 legitimate messages (public corpus, untrained Bayes).}
\label{tab:corpus}
\begin{tabular}{@{}lcc@{}}
\toprule
Threshold & spam detected & legitimate flagged (false positive) \\
\midrule
5.0 (configured) & 349/500 (69.8\%) & 52/500 (10.4\%) \\
3.0 & 438/500 (87.6\%) & 226/500 (45.2\%) \\
\bottomrule
\end{tabular}
\end{table}

% ------------------------------------------------------------------------------------------
\subsection{Route and signal hijacking: DNS hijacking}

\paragraph{Description.}
The attacker changes the answer to a name lookup so that clients reach a machine of the attacker's
choice instead of the real service. This work considers tampering with the authoritative zone, the case
behind the 2019 DNS infrastructure hijacking campaigns~\cite{cisa2019ed1901}.

\paragraph{Why it is relevant here.}
Clients reach the platform by name: \code{ingestion.ssle.local} and \code{analytics.ssle.local}
point at the node that publishes the NodePorts. Redirecting \code{ingestion} sends every sensor
reading to the attacker (data theft and silent data loss) and lets it return fake answers, while
every service stays up and no fault alarm fires.

\paragraph{How it can be exploited (PoC).}
The test overwrites the zone file \code{/etc/dnsmasq-zones/ssle.hosts} so that
\code{ingestion.ssle.local} points to an attacker container, while a client resolves the name every
0.25~s. A control phase repeats the attack with the DNS host's Wazuh agent stopped.

\paragraph{Monitoring and mitigation metrics.}
The DNS host's agent monitors \code{/etc/dnsmasq-zones} with real-time file-integrity monitoring
(all checks). A change raises Wazuh rule \rid{550} (integrity checksum changed), and the custom rule
\rid{100230} (level 12) matches it for the zone file. The client's resolution timeline measures the
exposure window.

\paragraph{Action-response.}
\code{restore-dns} runs on the DNS host: it copies the golden zone over the live one and restarts
dnsmasq, but only if they differ, so its own restore does not trigger it again.

\paragraph{Post-mortem.}
\rid{100230} fired 0.1~s after the hijack and the zone was restored 0.2~s after it on the k3s run
(0.3~s on Incus). The client never received the attacker's address (exposure 0.00~s at a 0.25~s
sampling interval). Without Wazuh the hijack was still in place after 15~s. Limitations: only
tampering with the zone \emph{file} is detected, not cache poisoning, a rogue DNS or DHCP server or
on-path spoofing, for which DNSSEC~\cite{rfc4033} would be needed; clients may keep a bad answer
for its TTL; root on the DNS host could change the golden copy or stop the agent. Inside the cluster
the equivalent attack would target the CoreDNS configuration, which file monitoring does not see;
it would need Kubernetes audit logs in Wazuh.
